%% ==========================================================================
%%  FINAL REPORT --- SECTION 2, THREAT MODEL AND SECURITY
%%  Source: docs/threat-model.md sections 1-3, 5, 7 and 8, and
%%  REQUIREMENTS.md section 3.2 for the ideal functionalities.
%% ==========================================================================
\section{Threat model and security}
\label{sec:threat}

This section says what \sys{} protects, against whom, and how far the argument
for it goes. It comes before the design on purpose: every later section is read
against it, and the most important sentences in this report are the ones that
say what is \emph{not} protected.

\subsection{Setting}
\label{sec:threat:setting}

There are three servers $\party{0}, \party{1}, \party{2}$, run by independent
operators, and $\users$ users. The rating matrix $\ratings \in \R^{\users \times
\items}$ is held in 2-of-3 replicated secret shares, so no single server holds
it. Training produces an item-embedding matrix $\itemembed \in \R^{\dime \times
\items}$, which is opened to everyone, and user embeddings $\userembed \in
\R^{\users \times \dime}$, which stay shared. Opening $\itemembed$ is
\nudgeplain{}'s design and not an oversight: each converged row is revealed
before the next is computed, which is what keeps orthogonalisation cheap.

To be served, user $i$ asks each server for its share of the score vector
$\avec{i}\itemembed$, reconstructs it, and selects the top $\topk$ items
\emph{on the user's own machine}. The user then fetches each of the $\topk$
records by two-server DPF-PIR against $\party{0}$ and $\party{1}$. The catalogue
of records is public, and all three servers hold it.

\subsection{Ideal functionalities}
\label{sec:threat:ideal}

Two functionalities describe what the system is meant to compute and to whom
each output goes. They are taken from the project's requirements document.
%
\begin{align*}
  \Ftrain(\shr{\ratings}, \dime, \iters)\colon\quad
    & (\userembed, \itemembed) \gets \ApproxFactor(\ratings) \text{ by $\iters$
      power iterations per component,}\\
    & \text{every server receives $\itemembed$ in the clear and its shares of
      $\shr{\userembed}$.}\\[4pt]
  \Fserve(\shr{\avec{i}}, \itemembed, \topk)\colon\quad
    & \mathbf{s} \gets \avec{i}\itemembed, \quad
      \mathbf{s}[j] \gets -\infty \text{ for every item $j$ user $i$ has
      rated},\\
    & T \gets \text{the top-$\topk$ indices of } \mathbf{s},\\
    & \text{the user receives $T$ and the records $D[T]$, and the servers
      receive nothing.}
\end{align*}
%
The second line of $\Fserve$'s output is what \sys{} adds to \nudgeplain{}. A
third functionality is needed because of one of our two normalisers, and
Section~\ref{sec:threat:gaps} explains why: $\Fleaky$ is $\Ftrain$ with the
sequence of norms $\{\norm{\mathbf{v}}\}$ from every normalisation added to the
servers' output.

\subsection{Adversary classes}
\label{sec:threat:classes}

Table~\ref{tab:adversaries} lists six classes. The headline guarantee is against
\textbf{A1}, one semi-honest server, which is exactly \nudgeplain{}'s model.

\begin{table}[t]
\centering
\small
\caption{Adversary classes and what \sys{} guarantees against each. A3 is the
assumption the whole system rests on.}
\label{tab:adversaries}
\begin{tabular}{@{}l p{3.6cm} p{9.2cm}@{}}
\toprule
 & Adversary & Guarantee \\
\midrule
A1 & One semi-honest server & Nothing is learned beyond the leakage profile of
  Section~\ref{sec:threat:leakage}. Argued by the sketch in
  Section~\ref{sec:threat:sketch}. \\
A2 & One semi-honest server and colluding users & Honest users' ratings stay
  hidden. Model \emph{quality} is not protected: colluding clients can distort
  the harvested training input (Section~\ref{sec:leak:malicious}). \\
A3 & Two or more colluding servers & \textbf{No guarantee.} Any two servers hold
  all three shares and reconstruct everything, and the two PIR servers together
  reconstruct every fetched index. \\
A4 & A malicious (deviating) server & Out of scope, as in \nudgeplain{}. A
  deviating server can corrupt correctness and availability. \\
A5 & A malicious user & Retrieval corrupts only that user's own output. The
  harvest path accepted a vote of any weight, which we close for inflation but
  not for redistribution, and full-domain evaluation is an unmitigated
  denial-of-service lever (Section~\ref{sec:leak:malicious}). \\
A6 & A network observer & Sees message sizes and timing. Sizes depend only on
  public parameters, measured in Section~\ref{sec:eval:wire}. Server-side timing
  is bounded in Section~\ref{sec:eval:timing}, and wire timing is not
  analysed. \\
\bottomrule
\end{tabular}
\end{table}

\subsection{Leakage profile}
\label{sec:threat:leakage}

\paragraph{Hidden from any single server.} Individual ratings, and which items a
user rated. The user embeddings $\userembed$. Each score vector, and the top-$k$
indices chosen from it. Which records a user fetches.

\paragraph{Leaked by design.} The parameters $\users, \items, \dime, \iters,
\topk, \bits, \fracbits$. The item-embedding matrix $\itemembed$, in the clear.
The number of non-zero ratings, which the truncation schedule uses as a public
bound and which \nudgeplain{} also leaks, unless every user is made to submit a
fixed number of ratings at a cost in quality. How many queries each user makes,
and when.

\paragraph{Outside the protocol's protection.} Correlation between when a user
queries and external events. Anything an adversary can infer from $\itemembed$
together with an observed fetch, which is precisely why the fetch has to be
private and is the subject of Section~\ref{sec:leak}.

\subsection{A real/ideal simulation sketch}
\label{sec:threat:sketch}

The claim, for class A1, is that for each corrupted server $\party{c}$ there is
a simulator $\Sim_c$ which, given only $\party{c}$'s input shares and the output
the functionality declares for it, produces a transcript computationally
indistinguishable from $\party{c}$'s real view. That view is its input shares,
its pseudorandom-function keys, its local randomness, and every ring element it
receives. We argue it step by step against the code. This is a \emph{sketch}.
There is no formal security parameter, no reduction, and no written hybrid
argument, and the three servers are not symmetric, since only $\party{0}$ and
$\party{1}$ hold DPF keys and the truncation protocol rotates a helper role. So
``by symmetry'' is never used.

\paragraph{The substrate.} A value is shared as $x = x_0 + x_1 + x_2$ with party
$i$ holding $(x_i, x_{i+1})$. Any two of three shares of a 3-out-of-3 additive
sharing are uniform and independent of $x$, so $\Sim_c$ samples two uniform ring
elements per value and never needs $\ratings$. Correlated randomness comes from
pairwise PRF keys. Party $i$ derives $\alpha_i = F(k_i, \mathit{ctr}) - F(k_{i-1},
\mathit{ctr})$, and the three sum to zero with no communication. $\party{c}$ can
compute its own $\alpha_c$ but no other party's, because each needs a key it
does not hold. That asymmetry is what makes every incoming element below
uniform.

\paragraph{Matrix--vector products.} Each party computes its share of the product
locally, masks it with $\alpha_i$, and sends one ring element per output entry.
$\party{c}$ receives $z_{c+1}$, masked by $\alpha_{c+1}$, which uses a key
$\party{c}$ lacks, so $\Sim_c$ samples it uniformly. The message size is set by
the public output length alone.

\paragraph{Truncation.} A helper $h$ supplies a correlated pair $(r, r \gg t)$
and the other two parties open $c = x - r$ and reshare. There are two cases. If
$\party{c}$ is an opener, it sees $x - r$ with $r$ uniform and known only to
$h$, so $\Sim_c$ samples $c$ uniformly. If $\party{c}$ is the helper, it sees
only the pair it generated, and the reshare is masked with a pairwise secret it
does not hold, so $\Sim_c$ replays its own randomness and samples the rest. A
party that held $r$ and also saw $x - r$ would recover $x$, which is why the
helper is passed explicitly at every call site.

\paragraph{Normalisation.} For the spec-faithful \texttt{FssNormalizer}, the
comparison gate opens $S + r$ where $r$ is drawn from a range $\kappa = 40$ bits
wider than $S$, so the opening is statistically close to uniform and $\Sim_c$
samples it. The DCF evaluations are local, and the Newton iterations are
multiplications and truncations already covered.

\paragraph{Orthogonalisation.} Entirely local, because every product is a share
times a public row of $\itemembed$. Nothing is sent.

\paragraph{Opening each row of $\itemembed$.} $\itemembed$ is $\Ftrain$'s
declared output, so $\Sim_c$ is given it and replays it. This step is free in the
simulation and is also the largest leakage surface in the system. That is not a
contradiction. It is leakage by design, and Section~\ref{sec:leak:attack} is the
measurement of what it costs.

\paragraph{Serving.} Scores are a public matrix applied to a shared vector, so
every server computes its share locally and sends it to the user, not to another
server. The seen-item mask arrives from the user as a share, which $\Sim_c$
samples.

\paragraph{Delivery.} $\party{c} \in \{\party{0}, \party{1}\}$ receives one DPF
key per fetch. A single key is pseudorandom and independent of the index by the
DPF's security, so $\Sim_c$ samples a key of the right length, 227 bytes at 11
domain bits. The server then expands the key over the whole domain and takes an
inner product with every record, so the work does not depend on the index.
$\party{2}$ receives nothing in this step.

\paragraph{Harvesting.} The difference $h_0 - h_1$ of the two PIR servers'
accumulated expansions is a 2-of-2 sharing of the consumption counts. It is
converted to 2-of-3 replicated form after re-randomising with a pairwise secret
that $\party{2}$ does not hold, so $\Sim_c$ samples uniformly. Without that
re-randomisation, $\party{2}$ would be sent $h_0$ in the clear alongside a zero
share and could subtract its way to the counts.

\subsection{Where the argument does not close}
\label{sec:threat:gaps}

\paragraph{1. The revealing normaliser is not simulable against $\Ftrain$.} The
cheaper normaliser opens $\norm{\mathbf{v}}^2$ at every normalisation, which is
$\dime(\iters + 1)$ scalars per run: the
trajectory of the power iteration's norms, which depend on the singular values of
$\ratings$. No simulator given only $\itemembed$ can produce them, because they
are not a function of $\itemembed$. We could exclude this path and claim
security only for the spec-faithful one. We instead claim security for it
against the weaker $\Fleaky$, because that is what we ran: \textbf{every
$\bits = 64$ quality number in this report comes from the revealing path}.
Claiming the strong statement while reporting numbers from the weak one would be
the dishonest version of this section. Section~\ref{sec:eval:quality} gives the
price of the path that does close.

\paragraph{2. Repeated training is argued, not proven.} Each round is covered
above, and harvesting adds no leak per round. What an adversary learns from many
successive openings of $\itemembed$, each trained on consumption that the
previous round's recommendations influenced, is not answered here.

\paragraph{3. Composition is sequential and assumed.} The protocols run one after
another with fresh correlated randomness, the standard setting for this kind of
argument, but no composition theorem is invoked and no universal-composability
claim is made.

\subsection{What we do not claim}
\label{sec:threat:notclaimed}

We claim nothing against two colluding servers, against a deviating server, or
against an observer who sees both PIR links at once, who reconstructs the record
by the definition of two-server PIR. We do not claim constant-time execution,
only a bound on index-dependent server time. We do not claim a proof. And we do
not claim that hiding the fetch hides the recommendation itself from anyone who
later sees what the user does with it.
